% 2de graad, Spreidingsdiagrammen (reeks Statistiek).
% Renoir: een menigte op een dansfeest, een hele groep in één beeld.
% Motief: een spreidingsdiagram met de trendlijn y = 0,648x + 0,418 (kleinste kwadraten, nagerekend).
\def\omslagbestand{omslag-galette.jpg}
\def\omslagbron{Pierre-Auguste Renoir, \textit{Bal du moulin de la Galette}, 1876 (publiek domein)}
\def\omslagrand{10.6 + 0.35*sin(\x*360/21)}
\def\omslagmotief{%
  \randlijn
  \begin{scope}[shift={(15,4.6)}]
  \draw[white,line width=.7pt,-{Stealth[length=2mm]}] (-.2,0) -- (4.8,0);
  \draw[white,line width=.7pt,-{Stealth[length=2mm]}] (0,-.2) -- (0,4);
  \draw[fel,line width=1.6pt] (0,.418) -- (4.5,3.334);
  \foreach \x/\y in {0.5/0.6,1/1.4,1.5/1.2,2/1.9,2.5/1.7,3/2.6,3.5/2.4,4/3.2} \fill[white] (\x,\y) circle (1mm);
  \end{scope}
}
